\documentclass[11pt,a4paper]{article}
\pdfoutput=1

\usepackage{geometry}                % See geometry.pdf to learn the layout options. There are lots.
\geometry{letterpaper}                   % ... or a4paper or a5paper or ... 
%\geometry{landscape}                % Activate for for rotated page geometry
%\usepackage[parfill]{parskip}    % Activate to begin paragraphs with an empty line rather than an indent
\usepackage{graphicx}
\usepackage{amssymb}
%\usepackage{epstopdf}
%\DeclareGraphicsRule{.tif}{png}{.png}{`convert #1 `dirname #1`/`basename #1 .tif`.png}
\usepackage{xcolor}
\usepackage{amsmath}
\usepackage{subfigure}
\usepackage{listings}    
\usepackage{bbm}
\usepackage{cancel}
\usepackage{relsize}
%\usepackage{listings}
\usepackage{xcolor}
%\usepackage{tcolorbox}
\usepackage[most]{tcolorbox}
\usepackage{hyperref}
\usepackage{mathtools}

\hypersetup{
    colorlinks=true,
    linkcolor=blue,
    filecolor=magenta,      
    urlcolor=blue,
    pdftitle={Overleaf Example},
    pdfpagemode=FullScreen,
    }

\definecolor{codegreen}{rgb}{0,0.6,0}
\definecolor{codegray}{rgb}{0.5,0.5,0.5}
\definecolor{codepurple}{rgb}{0.58,0,0.82}
\definecolor{backcolour}{rgb}{0.95,0.95,0.92}

\pagecolor{white}

%%%%%%% Borun's commands %%%%%%%%%%%%%%

\newcommand\be{\begin{equation}}
\newcommand\bea{\begin{eqnarray}}
\newcommand\ee{\end{equation}}
\newcommand\eea{\end{eqnarray}}
\newcommand\Regge{\alpha'}
\newcommand{\bdm}{\begin{displaymath}}
\newcommand{\edm}{\end{displaymath}}
\newcommand{\nn}{\nonumber \\}
\newcommand{\f}[2]{\frac{#1}{#2}}
\newcommand{\bref}[1]{(\ref{#1})}
\newcommand\h{\frac{1}{2}}
\newcommand{\ket}[1]{|#1 \rangle}
\newcommand{\bra}[1]{\langle #1 |}
\newcommand{\e}[2]{\mathbb E_{#2}[#1]}
\newcommand{\todo}[1]{{\color{cyan}{TODO: #1}}}

\newcommand{\partitionfunctionintegrand}{{\color{olive} \Bigg(  \prod_{t=0}^{T-1} p(S_{t+1}|S_t,A_t) \pi(A_t|S_t) \Bigg) p(S_0) }}
\newcommand{\partitionfunctionintegrandtheta}{{\color{olive} \Bigg(  \prod_{t=0}^{T-1} p(S_{t+1}|S_t,A_t) \pi_\theta(A_t|S_t) \Bigg) p(S_0) }}
\newcommand{\partitionfunctionintegrandBehaviourial}{{\color{olive} \Bigg(  \prod_{t=0}^{T-1} p(S_{t+1}|S_t,A_t) b(A_t|S_t) \Bigg) p(S_0) }}
\newcommand{\partitionfunctionintegral}{{\color{olive} \sum_{S_{0:T}, A_{0:T-1}}}}

\newcommand{\partitionfunctionintegrandsecond}{{\color{olive} \Bigg(  \prod_{t=0}^{T-1} p(S_{t+1}|S_t,A_t) \pi(A_t|S_t) \Bigg) }}
\newcommand{\partitionfunctionintegralsecond}{{\color{olive} \sum_{S_{1:T}, A_{0:T-1}}}}

\newcommand{\argmax}[1]{\underset{#1}{\text{argmax }}}

\title{Reinforcement Learning}
\author{Borun D. Chowdhury}
%\date{}                                           % Activate to display a given date or no date

\begin{document}
\maketitle
\tableofcontents

\section{Introduction}

Imagine a problem where we have an {\em agent} that interacts with an {\em environment} and {\em we} want it to achieve a goal without explicitly telling it what to do? 

We could assign a {\em numerical return} to the outcomes with the largest value for the goal would like to agent to achieve and simply tweak the agent to take {\em actions} (out of the set of possible actions at any given time) so as to maximize the return.  

The return itself could come from rewards sprinkled along the way (think stock market returns) or from a final reward at the last step (checkmate!). In practice, given a goal, one has to decide on numerical rewards and compose a final return from those rewards that is consistent with the goal.

Let's make this concrete. Then a sequence of states and actions are written as 
\be
(S_0,A_0,\cancel{R_0}), (S_1,A_1,R_1),(S_2,A_2,R_2), \dots (S_T,\cancel{A_T},R_T) 
\ee
Here, we start in a state $S_0$ (which represents the {\em state of the environment as well as the internal state of the agent}) and may be only partially visible to the agent. The agent then takes the action $A_0$. This results in reward which we {\em choose to label $R_1$}. This explains why $R_0$ is cancelled out since the game has just started. This goes on till the game terminates at time $T$. Note it is not necessary that a game terminates.

This process is often made more complicated in practice because of inherent stochasticity in the transition of states after an action is taken as well as in the rewards generated. So in the notation above
\begin{itemize}
\item The game may start in a state $S_0 \sim p(S_0)$
\item After the agent taken the action $A_t$ the next state may be sample $S_{t+1} \sim p(S_{t+1}|S_t,A_t)$
\item The reward is sampled $R_{t+1} \sim p(R_{t+1}|S_t,A_t)$
\end{itemize}

Note we have given in to the Markovian fantasy above in that we assume the next step depends only on the current state. In theory, this is not a big assumption, as one can always {\em define} the state to contain all the past information that is required. In practice though it is a strong assumption that is frequently made without much thought largely due to the complexity involved in the alternative assumption.

Now with this Markovian fantasy the thing the agent needs to do at time $t$ is to decide what action to take in the state $S_t$. In order to do this the agent needs to have a {\em policy} from which actions are sampled $A_t \sim \pi(A_t|S_t)$. This may be a deterministic or stochastic policy. Now we come to the two main problems of Reinforcement Learning
\begin{itemize}
\item What are $p(R_{t+1} | S_t,A_t)$ and $p(S_{t+1}|S_t,A_t)$?
\item What should $\pi(A_t|S_t)$ be?
\end{itemize}

The rest of the lecture will be  discussing how to solve these problems but we briefly discuss the general ideas. First we need to get back to discussing the {\em return} - a numerical proxy for the goal instead of returns received immediately after an action is taken. The return would be some aggregating function of the rewards in the future of an action and for most problems the contribution  of temporally distant reward ought to be smaller. For computational reason we take
\be
G_t  = \sum_{k=0}^\infty \gamma^k R_{t+k+1}
\ee

Given a policy $\pi(A_t|S_t)$ we can {\em in theory} figure out $p(R_{t+1} | S_t,A_t)$ and $p(S_{t+1}|S_t,A_t)$ by simply sampling these repeatedly. Note these would be independent of the policy $\pi$ as long as the policy leads to visiting the concerned state and take the corresponding action. We would also that way end up getting $p(G_t|S_t,A_t;\pi)$. This is the Monte Carlo method. For details see section~\ref{sec:monte_carlo}.

There is also the bootstrapping method where we make initial guesses for $q_\pi(s,a)$ and take a few steps and update the initial guess by the reward obtained in the process and the best guess at the time. For details see section~\ref{sec:bootstrapping_methods}. 


Once we have found the returns for a policy, we can then modify the policy to improve the returns. This is called policy improvement. 

In practice one doesn't go all the way to complete the evaluation for a given policy before improving it. All the flavours of going back and forth between evaluation and improvement are termed {\em Generalized Policy Iteration}.



\section{Markov Decision Processes}

A MDP is a sequence of states, actions and rewards 
\be
(S_0,A_0,\cancel{R_0})),(S_1,A_1,R_1),(S_2,A_2,R_2) \dots (S_T,\cancel{A_T},R_T) \label{MDP}
\ee
that may or may not have a terminal state. If it does have a terminal state then we call the full run an {\em episode}.

The {\em process} part should be obvious but let's focus on the other terms. First let's discuss the term {\em Markov}. We make the  Markovian assumptions that we can capture the process by just modelling $p(S_{t+1}|S_t,A_t)$ and $\pi(A_t|S_t)$. The former is the probability to transfer to a new state given a state and an action while  the latter is the probability to take an action $A_t$ given that the state is $S_t$. Strictly speaking, this is not an assumption as given a big enough {\em state space}, we can always capture the past states that matter in the current state. However, in practice this becomes an approximation because the efforts to enlarge the state space is based on feature engineering. Next, let's discuss the term {\em decision}. The term $\pi(A_t|S_t)$ is taken to be {\em tunable} and in that sense it is a decision being taken.

 %Note that the 
 %\be
 %p(S_{t+1},A_{t+1}| S_t,A_t) = \pi(A_{t+1}|S_{t+1}) p(S_{t+1}|S_t,A_t)
 %\ee
 
For an episodic MDP we can write the partition function
\bea
Z_\pi = \partitionfunctionintegral \partitionfunctionintegrand
\eea
The continuing case is the formal $T \to \infty$ limit of the above.
 
This gives
\bea
\mathbb E_\pi[R_{\tau+1}] &=&  \partitionfunctionintegral R_{\tau+1} p(R_{\tau+1}|S_\tau,A_\tau) \partitionfunctionintegrand\nn
&=& {\color{brown} \sum_{S_{0:\tau}, A_{0:\tau}}}  R_{\tau+1} p(R_{\tau+1}|S_\tau,A_\tau) {\color{brown} \prod_{t=0}^{\tau} p(S_{t+1}|S_t,A_t) \pi(A_t|S_t)  p(S_0)}
\eea
The actions taken and states coming  after the reward in question do not impact the expectation value.  However, formally we will always take the product over all times as it doesn't make a difference.  

So if the agent is about to decide which action to take at time $t$ they should only care about $\mathbb E[R_{t+k}]~\forall~k\ \ge 1$. We can then say that the agent ought to decide based on
\be
\sum_{k \ge 1} \mathbb E_\pi[R_{t+k}|S_t,A_t]
\ee
however if $k$ is unbounded then this function is unbounded. We thus define a stochastic reward 
\bea
G_t &\coloneqq& \sum_{k=1}^\infty \gamma^{k-1} R_{t+k} \label{def:G} \nn
&=& R_{t+1} + \gamma G_{t+1}
\eea
and the agent then takes actions based on
\bea
A _t= \argmax{a} \mathbb E_\pi[G_t|S_t,a]
\eea

\section{Value Functions}

We define the {\em value function}
\bea
v_\pi(s) &\coloneqq& \mathbb E_\pi[G_0|S_0=s] \nn
&=& \sum_{t=0}^\infty \gamma^t \mathbb E_\pi[R_{t+1} |S_0=s]
\eea
and the {\em state-action value function}
\bea
q_\pi(s,a) &\coloneqq& \mathbb E_\pi[G_0|S_0=s,A_0=a] \nn
&=& \sum_{t=0}^\infty \gamma^t \mathbb E_\pi[R_{t+1} |S_0=s,A_0=a]
\eea


Since the process is Markovian the transition probabilities are independent of the time step. Below we define
\bea
p(s'|s,a) \coloneqq p(S_{t+1}=s'|S_t=s,A_t=a)
\eea
and likewise other expressions.

We use
\be
\mathbb E[A] =\mathbb E[ \mathbb E[A|B] ]
\ee
extensively.

We then get
\bea
v_\pi(s) &=& \e{G_t|S_t=s}{\pi} \nn
&=& \e{R_{t+1} + \gamma G_{t+1} |S_t=s}{\pi} \nn
&=& \e{R_{t+1}|S_t=s}{\pi} + \gamma \e{G_{t+1}|S_t=s}{\pi} \nn
&=& \e{R_{t+1}|S_t=s}{\pi} + \gamma \e{\e{G_{t+1}|S_{t+1}=s'}{\pi}|S_t=s}{\pi}  \nn
&=& \e{R_{t+1}|S_t=s}{\pi} + \gamma \e{v_\pi(s') |S_t=s}{\pi} \nn
&=& \sum_a \Bigg[\e{R_{t+1}|S_t=s,a}{\pi} + \gamma \e{v_\pi(s') |S_t=s,a}{\pi} \Bigg] \pi(a|s) \nn
&=& \sum_a \Bigg[\sum_r r p(r|s,a) + \gamma \sum_{s'} v_\pi(s') p(s'|s,a) \Bigg] \pi(a|s) 
\eea



Similarly, we get
\bea
q_\pi(s,a) &=& \e{G_t | S_t=s,A_t=a}{\pi} \nn
&=& \e{R_{t+1} |S_t=s,A_t=a}{\pi} + \gamma \e{G_{t+1} | S_t=s, A_t=a}{\pi} \nn
&=& \e{R_{t+1} |S_t=s,A_t=a}{\pi} + \gamma \e{ \e{G_{t+1} | S_{t+1}=s',A_{t+1}=a'}{\pi}|S_t=s, A_t=a}{\pi} \nn
&=& \e{R_{t+1} |S_t=s,A_t=a}{\pi} + \gamma \e{ q_\pi(s',a')|S_t=s, A_t=a}{\pi} \nn
&=& \sum_r r p(r|s,a) + \gamma \sum_{s',a'} q_\pi(s',a') p(s',a'|s,a) \nn
&=&  \sum_r r p(r|s,a) + \gamma \sum_{s',a'} q_\pi(s',a') \pi(a'|s') p(s'|s,a)
\eea

From these we see
\bea
v_\pi(s) &=&  \e{G_t|S_t=s}{\pi} \nn
&=& \e{\e{G_t|S_t=s,A_t=a}{\pi} }{\pi} \nn
&=& \e{q_\pi(s,a)}{\pi} \nn
&=& \sum_a q_\pi(s,a) \pi(a|s)
\eea
Further we have
\bea
q_\pi(s,a) = \sum_r r p(r|s,a) + \gamma \sum_{s'} v_\pi(s') p(s'|s,a)
\eea


If these {\em were} to be known accurately, the agent can simply take action $A_t = \argmax{a} q_\pi(S_t,a)$. 

However, easier said than done, there are two problems to be solved
\begin{itemize}
\item Actually find the correct values for $q_\pi(s,a)$
\item Find the best $\pi_\star$ such that $\max_a q_{\pi_*}(s,a) \ge \max_a q_\pi(s,a)~~ \forall ~\pi,s$
\end{itemize}
We get for the optimal policy
\bea
v_\star(s) &=&\underset{a}{\text{max}}  \Bigg[\sum_r r p(r|s,a) + \gamma \sum_{s'} v_\star(s') p(s'|s,a) \Bigg]  \\
q_\star(s,a) &=& \sum_r r p(r|s,a)  + \gamma \sum_{s'} \left( \underset{a'}{\text{max}} ~q_\star(s',a') \right)p(s'|s,a)
\eea






\section{Monte Carlo} \label{sec:monte_carlo}

The expected reward is
\bea
\mathbb E_\pi[G_0|S_0] &=& \partitionfunctionintegralsecond  \Bigg( \sum_{R_{1:T}} \sum_{\tau=0}^{T-1} \gamma^\tau R_{\tau+1} p(R_{\tau+1}| S_\tau,A_\tau) \Bigg) \partitionfunctionintegrandsecond  \nn
\eea
If we repeatedly sample the MDP then the average value will give us the above on expectation. Furthermore, because it's a Markovian process, there is nothing special about the starting state so the multiple samplings of MDP will give us 
\bea
q_\pi(s,a) = \mathbb E[G_0|S_0 = s,A_0=a]~\forall~s,a
\eea
as long as the $(s,a)$ pair are sampled sufficiently.


A very useful idea in Monte Carlo is that of {\em importance sampling}. Here, we may sample using policy $b$ while learning the expectations for policy $\pi$. 

Concretely, if we have two different policies $\pi$ and $b$ 
\bea
E_\pi[R_{\tau+1}]  = \partitionfunctionintegral  R_{\tau+1} p(R_{\tau+1}|S_\tau,A_\tau) \partitionfunctionintegrand \nn
E_b[R_{\tau+1}]  = \partitionfunctionintegral R_{\tau+1} p(R_{\tau+1}|S_\tau,A_\tau) \partitionfunctionintegrandBehaviourial \nn
\eea
Then we get 
\bea
E_\pi[R_{\tau+1}] =\mathbb E_{b}[ \left( \prod_{t=0}^{\tau} \frac{\pi(A_{t} | S_{t} )}{b(A_{t} | S_{t} )} \right)   R_{\tau+1}] 
\eea

So we get
\bea
\mathbb  E_\pi [ G_0|S_0] &=& \sum_{\tau=0}^{T-1} \gamma^\tau \mathbb E_{b}[ \left( \prod_{t=0}^{\tau} \frac{\pi(A_{t} | S_{t} )}{b(A_{t} | S_{t} )} \right)   R_{\tau+1}]  \nn
&=& \partitionfunctionintegralsecond  \Bigg( \sum_{R_{1:T}} \sum_{\tau=0}^{T-1} \gamma^\tau R_{\tau+1} \left( \prod_{t=0}^{\tau} \frac{\pi(A_{t} | S_{t} )}{b(A_{t} | S_{t} )} \right) p(R_{\tau+1}|S_\tau,A_\tau) \Bigg)  \nn
&\times& \partitionfunctionintegrandsecond \nn
\eea
Note this is called {\em per-decision} importance sampling. Another version, one without a name, would have each $R_t$ taken with $\prod_{n=0}^{T-1} \frac{\pi(A_n|S_n)}{b(A_n|S_n)}$. Since the actions after a reward do not matter, the extra terms just add to the variance.

\section{Policy Gradient}

We define
\bea
J_\theta &=& \partitionfunctionintegral  \Bigg( \sum_{R_{1:T}} \sum_{\tau=0}^{T-1} \gamma^\tau R_{\tau+1} p(R_{\tau+1}| S_\tau,A_\tau) \Bigg) \partitionfunctionintegrandtheta  \nn
\eea
We then have
\bea
\partial_\theta J_\theta &=& \partitionfunctionintegral  \Bigg( \sum_{R_{1:T}} \sum_{\tau=0}^{T-1} \gamma^\tau R_{\tau+1} p(R_{\tau+1}| S_\tau,A_\tau)  \Bigg)  \nn
&\times& \Bigg( \sum_{\eta=0}^{T-1} \partial_\theta \log \pi(A_\eta |S_\eta) \Bigg) \partitionfunctionintegrandtheta \nn
&=& \partitionfunctionintegral  \Bigg( \sum_{\eta=0}^{T-1}  \sum_{R_{1:T}} \sum_{\tau=0}^{T-1} \gamma^\tau R_{\tau+1}  (\partial_\theta \log \pi(A_\eta |S_\eta) ) p(R_{\tau+1}| S_\tau,A_\tau)  \Bigg)  \nn
&\times&  \partitionfunctionintegrandtheta \nn
&=& \partitionfunctionintegral  \Bigg( \sum_{\eta=0}^{T-1} \partial_\theta \log \pi(A_\eta |S_\eta) \gamma^\eta \sum_{R_{\eta+1:T}} \sum_{k=0}^{T-1-\eta}  \gamma^k R_{\eta+k+1}   p(R_{\eta+k+1}| S_{\eta+k},A_{\eta+k})  \Bigg)  \nn
&\times&  \partitionfunctionintegrandtheta 
 \eea
 We can use the relation between returns and rewards to rewrite the above as
 \bea
  \partial_\theta J_\theta &=& \partitionfunctionintegral  \sum_{R_{1:T}}   \Bigg( \sum_{\eta=0}^{T-1} \gamma^\eta G_\eta \partial_\theta \log \pi(A_\eta|S_\eta) \Bigg) {\color{olive} \prod_{t=0}^{T-1} } p(R_{t+1}|S_t,A_t) {\color{olive}p(S_{t+1}|S_t,A_t) \pi_\theta(A_t|S_t) p(S_0)} \nn
\eea





\section{Bootstrapping Methods} \label{sec:bootstrapping_methods}

Bootstrap methods rely on {\em assuming} some initial value and updating it by partial traversals of the MDP.  In general they are of the form

\bea
Q_{t+n}(S_t,A_t) =  Q_{t+n-1} (S_t,A_t) +  \alpha ~ \text{something}
\eea
where $\alpha$ is a hyperparameter controlling the EMA scale of updates and the extra part is the actual update. This is the part that will differ in the various methods described below.

The result would converge when
\be
\e{\text{something}}{b} = 0
\ee
where $b$ is the behaviour policy. The form of the correction would determine what the Q-values actually converge to.



\subsection{On Policy aka  SARSA}

Let us define a generalized or partial return as
\bea
G_{t:t+n} \coloneqq \sum_{k=1}^{n} \gamma^{k-1} R_{t+k} + \gamma^n Q_{t+n-1}(S_{t+n},A_{t+n})
\eea
We see this is a natural generalisation of equation~\ref{def:G} where we take the actual returns for $n$ steps and then the best estimate {\em at the time} of the state-action value of the resulting state after those $n$ steps and the action taken thereafter. 

Also note the recursive relation
\bea
G_{t:t+n} = R_{t+1} + \gamma G_{t+1:t+n} 
\eea
and the boundary condition
\bea
G_{t+n:t+n} = Q_{t+n-1}(S_{t+n},A_{t+n})
\eea

The update rule is then
\bea
Q_{t+n} (S_t,A_t)  &=&  Q_{t+n-1}(S_t,A_t) + \alpha (G_{t:t+n} - Q_{t+n-1}(S_t,A_t))
\eea
and the Q-values learnt are on-policy.

This $n=1$ version of this is called SARSA and has the simple form
\bea
Q_{t+1} (S_t,A_t)  &=& Q_{t}(S_t,A_t) + \alpha [R_{t+1} + \gamma Q_t(S_{t+1},A_{t+1}) - Q_{t}(S_t,A_t)]
\eea
The general one can then be called the $n$ step SARSA and writing more explicitly we get
\bea
Q_{t+n} (S_t,A_t)  &=& Q_{t+n-1}(S_t,A_t) + \alpha [\sum_{k=1}^{n} \gamma^{k-1} R_{t+k} + \gamma^n Q_{t+n-1}(S_{t+n},A_{t+n})- Q_{t+n-1}(S_t,A_t) ] \nn
\eea

\subsection{Off Policy aka Q-Learning}

We us define
\be
\rho_{t:t+n} = \prod_{\tau=t}^{min(t+n,T-1)} \frac{\pi(A_\tau|S_\tau)}{b(A_\tau|S_\tau)}
\ee
where $\pi$ is the target policy and $b$ is the behaviour policy.

We define the partial returns for off policy as 
\bea
\tilde G_{t:t+n} = R_{t+1} + \gamma \rho_{t+1:t+1} \tilde G_{t+1:t+n}
\eea
with the boundary condition
\be
\tilde G_{t+n:t+n} = Q_{t+n-1}(S_{t+n},A_{t+n})
\ee

The the $n$ step SARSA can be replaced by
\bea
Q_{t+n} (S_t,A_t) &=& Q_{t+n-1}(S_t,A_t)  + \alpha [ \tilde G_{t:t+n} - \rho_{t+1:t+n-1} Q_{t+n-1}(S_t,A_t)] \nn
&=& Q_{t+n-1}(S_t,A_t) \nn
&+& \alpha \Bigg [\sum_{k=1}^{n} \gamma^{k-1}  \rho_{t+1:t+k-1} R_{t+k}  +  \rho_{t+1:t+n}  \gamma^n Q_{t+n-1}(S_{t+n},A_{t+n})  \nn
&-&  \rho_{t+1:t+n-1} Q_{t+n-1}(S_t,A_t) \Bigg]
\eea
when we sample by policy $b$ but want expectations as per $\pi$ as the policy.  Notice how we needing an importance sampling factor in front of $Q_{t+n-1}(S_t,A_t)$ in the TD error. This ensures that at convergence we have
\bea
&&\e{ \sum_{k=1}^{n} \gamma^{k-1}  \rho_{t+1:t+k-1} R_{t+k}  +  \rho_{t+1:t+n}  \gamma^n Q_{t+n-1}(S_{t+n},A_{t+n})  -  \rho_{t+1:t+n-1} Q_{t+n-1}(S_t,A_t) }{b} \nn
&=& \e{ \sum_{k=1}^{n} \gamma^{k-1} R_{t+k}  +  \gamma^n Q_{t+n-1}(S_{t+n},A_{t+n})  -  Q_{t+n-1}(S_t,A_t) }{\pi} \nn
&=&0
\eea
which means that the Q-values being learnt are those of the target policy $\pi$.




Now notice that if we don't mind extra variance then we can instead write this as
\bea
Q_{t+n} (S_t,A_t) &=&  Q_{t+n-1}(S_t,A_t) + \alpha \rho_{t+1:t+n}  [\sum_{k=0}^{n-1} \gamma^k R_{t+1+k} + \gamma^n Q_{t+n-1}(S_{t+n},A_{t+n}) - Q_{t+n-1}(S_t,A_t) ] \nn
\eea



Now let's look at the $n=1$ case. We get
\bea
Q_{t+1} (S_t,A_t) &=& Q_{t}(S_t,A_t) + \alpha [  R_{t+1}  + \gamma  \frac{\pi(A_{t+1}|S_{t+1})}{b(A_{t+1}|S_{t+1})}  Q_{t}(S_{t+1},A_{t+1}) - Q_{t}(S_t,A_t) ] \nn
\eea
If $\pi$ is the greedy policy we get have $\pi(A_{t+1}|S_{t+1}) =\delta_{\argmax{a}Q_t(S_{t+1},a),A_{t+1}}$. This gives 
\bea
Q_{t+1} (S_t,A_t) &=& Q_{t}(S_t,A_t) + \alpha [  R_{t+1}  + \gamma Q_{t}(S_{t+1},A_{t+1})  \frac{\delta_{\argmax{a} Q_t(S_{t+1},a),A_{t+1}}}{b(A_{t+1}|S_{t+1})}  - Q_{t} (S_t,A_t)] \nn
\eea
Furthermore, since
\bea
&&\mathbb E_b[ Q_t(S_{t+1},A_{t+1} ) \delta_{\argmax{a} Q_t(S_{t+1},a),A_{t+1}}\frac{\pi(A_{t+1}|S_{t+1})}{b(A_{t+1}|S_{t+1})} ]  \nn
&=& \mathbb E_\pi [ Q_t(S_{t+1},A_{t+1} )\delta_{\argmax{a} Q_t(S_{t+1},a),A_{t+1}}] \nn
&=& \max_a Q_t(S_{t+1},a)
\eea
we can replace the stochastic term with it's average 
\bea
Q_{t+1} (S_t,A_t) &=& Q_{t}(S_t,A_t) + \alpha [  R_{t+1}  + \gamma \max_a Q_t(S_{t+1},a) - Q_{t}(S_t,A_t)] 
\eea
This has reduced variance than the sampled form of the expression.



\subsection{Tree Backup}

For this the recursive relation for the partial return is modified by the red and blue pieces below
\bea
\tilde G_{t:t+n}  &=& R_{t+1} + \gamma \left( {\color{red} \sum_{a \ne A_{t+1}} Q_{t+n-1}(S_{t+1},a) \pi(a|S_{t+1}) } + {\color{blue} \pi(A_{t+1}|S_{t+1}) } \tilde G_{t+1:t+n} \right) \nn
\eea
Thus, while in Monte Carlo we put the full weight on the path taken by taking the red part as zero and the blue part as one, here we take the values of paths not taken into account as well.

Let's write down a couple of these 

\bea
\tilde G_{t:t+1}  = R_{t+1} + \gamma \sum_a \pi(a|S_{t+1}) Q_{t}(S_{t+1},a)
\eea
and
\bea
\tilde G_{t:t+2} &=&  R_{t+1} + \gamma \sum_{a \ne A_{t+1}} \pi(a|S_{t+1}) Q_{t+1} (S_{t+1},a) \nn
&& + \gamma \pi(A_{t+1}|S_{t+1})  \left(R_{t+2} + \gamma \sum_a \pi(a|S_{t+2}) Q_{t+1}(S_{t+2},a)
 \right)
\eea

The update rule is
\bea
Q_{t+n}(S_t,A_t) =Q_{t+n-1}(S_t,A_t)  + \alpha [ \tilde  G_{t:t+n} - Q_{t+n-1} (S_t,A_t)]
\eea
and the Q value learnt is as per the target policy.

\appendix
\section{Some results about Importance Sampling}

\bea
\mathbb E_b [ f(A_t) \frac{\pi(A_t|S_t)}{b(A_t|S_t)}] &=& \sum_a f(a) \frac{\pi(a|S_t)}{b(a|S_t)} b(a|S_t) \nn
&=& \sum_a f(a) \pi(a|S_t) \nn
&=& \mathbb E_\pi[ f(A_t)]
\eea

Therefore it is easy to see that
\be
\mathbb E_b[ c  \frac{\pi(A_t|S_t)}{b(A_t|S_t)}] = c
\ee
where $c$ is not a function of $A_t$.

\section{Bandits}

\subsection{Value Based Models}

Suppose there is a process where one can choose among a set of actions $\{a\}$ that give rewards $R$ then it makes sense to take actions that maximize the reward
\bea
q_*(a) = \e{R|A=a}{\pi_*}
\eea
which is the expectation of the reward given the action chose is $a$.

We want to 
\begin{itemize}
\item Find the true values for each action
\be
Q_t(a) := \frac{\sum_{i=1}^{t-1} R_i \delta_{A_i,a}}{\sum_{i=1}^{t-1}  \delta_{A_i,a}} \label{rewardrunningmean}
\ee
\item Take the action with the max value
\be
A_t := \underset{a}{\text{argmax}} ~ Q_t(a)
\ee
\end{itemize}

However, the tricky part is to do them together. For this we can 
\begin{itemize}
\item Take an $\epsilon$-greedy approach where we exploit with probability $1-\epsilon$ and explore with probability $\epsilon$,
\item Use the upper-confidence-bound method (\todo{find the math behind this})
\be
A_t := \underset{a}{\text{argmax}} ~ \left[ Q_t(a) + c \sqrt{\frac{\ln(t)}{N_t(a)}}\right]
\ee
\end{itemize}

The computation of the running mean for rewards eqn \ref{rewardrunningmean} per action naively requires keeping track of all rewards but we can instead do
\be
Q_{n+1} = Q_n + \frac{1}{n} [R_n - Q_n]
\ee
Furthermore if the problem is non-stationary we can instead have a fixed parameter $\alpha$ to exponential weight the prior reward
\bea
Q_{n+1} &=& Q_n + \alpha [R_n - Q_n] \nn
&=& (1-\alpha) Q_n + \alpha R_n \nn
&=& \alpha R_n + (1-\alpha) \alpha R_{n-1} + (1-\alpha)^2 \alpha R_{n-2}  + \dots \nn
&=& (1-\alpha)^n Q_1 + \alpha \sum_{i=1}^n (1-\alpha)^{n-i} R_i
\eea

\section{Monte Carlo}

\subsection{The idea}

When we do not know the dynamics of the system (i.e. $p(s'|s,a)$ and/or $p(r|s,a)$ then we can sample whole episodes several times to get
\be
v_\pi(s), q_\pi(s,a)
\ee 
by starting off in state $s$ or in state $s$ and take action $a$ and then follow through till the end of the episode. In doing so we can take the first visit approach where we update the averages for a state or state action pair only for the first visit or we can do this for all the visits. Both converge but proving theorems is easier for the former.


\section{Taking Averages For Monte Carlo Reinforcement Learning}

\subsection{Ordinary Importance Sampling}

For Monte Carlo we can talk about two kinds of averages across episodes. The first is {\em ordinary importance sampling} 
\bea
V(s) = \frac{\sum_{t \in \mathcal T(s)} \rho_{t:T(t)-1} G_t}{|\mathcal T(s)|}
\eea
The way to interpret this is that $\mathcal T(s)$ is the union of all time steps at which the state was $s$.  If we are discussing first visit MC then this would be restricted to first visit. T(t) is the first time of termination following $t$. $G_t$ are returns pertaining to time step $t$ and $ \rho_{t:T(t)-1}$ are the importance sampling ratios.

To show how this works we consider an example. We show how the values are updated for first visit and every visit MC for the first two episodes. 
\begin{itemize}
\item Episode 1: $\{(s,a),(s',a',R_1),(s,a,R_2),(s_t,R_3\}$ 
We have
\bea
G_0 &=& R_1 + \gamma R_2 + \gamma^2 R_3 \\
G_1 &=&  R_2 + \gamma R_3 \\
G_2 &=& R_3 
\eea

\begin{itemize}
\item First Visit
\bea
q(s,a) &=& G_0 \rho_{1,\infty} \nn
&=&  (R_1 + \gamma R_2+\gamma^2 R_3) \frac{\pi(a'|s') \pi(a|s)}{b(a'|s') b(a|s)} \\
q(s',a') &=& G_1 \rho_{2,\infty}  \nn
&=& (R_2 + \gamma R_3) \frac{\pi(a|s)}{b(a|s)}
\eea
\item Every Visit
\bea
q(s,a) &=& \frac{G_2 + G_0 \rho_{1,\infty}}{2} \nn
&=&  \frac{(R_1 + \gamma R_2+\gamma^2 R_3) \frac{\pi(a'|s') \pi(a|s)}{b(a'|s') b(a|s)}+R_3}{2} \\
q(s',a') &=& G_1 \rho_{2,\infty}  \nn
&=& (R_2 + \gamma R_3) \frac{\pi(a|s)}{b(a|s)}
\eea
\end{itemize}
\item  Episode 2: $\{(s',a'),(s',a',R_4),(s,\bar a,R_5),(s_t,R_6)\}$
\bea
G_0 &=& R_4 + \gamma R_5 + \gamma^2 R_6 \\
G_1 &=& R_5 + \gamma R_6 \\
G_2 &=& R_6
\eea
\begin{itemize}
\item First Visit
\bea
q(s,a) &=& (R_1 + \gamma R_2+\gamma^2 R_3) \frac{\pi(a'|s') \pi(a|s)}{b(a'|s') b(a|s)} \\
q(s',a') &=& \frac{(R_2 + \gamma R_3) \frac{\pi(a|s)}{b(a|s)} + (R_4 + \gamma R_5 + \gamma^2 R_6) \frac{\pi(a'|s') \pi(\bar a|s)}{b(a'|s') b(\bar a|s)}}{2} \\
q(s,\bar a) &=& R_6
\eea
\item Every Visit
\bea
q(s,a) &=& \frac{(R_1 + \gamma R_2+\gamma^2 R_3) \frac{\pi(a'|s') \pi(a|s)}{b(a'|s') b(a|s)}+R_3}{2} \\
q(s',a') &=& \frac{(R_2 + \gamma R_3) \frac{\pi(a|s)}{b(a|s)} + (R_4 + \gamma R_5 + \gamma^2 R_6) \frac{\pi(a'|s') \pi(\bar a|s)}{b(a'|s') b(\bar a|s)} + (R_5+ \gamma R_6)  \frac{ \pi(\bar a|s)}{ b(\bar a|s)}}{3} \\
q(s,\bar a) &=& R_6
\eea
\end{itemize}
\end{itemize}

\subsection{Ordinary Importance Sampling: Per-importance importance sampling}

At this point it makes sense to {\em re-write the above after dropping the importance sampling ratios after the reward} as they do not make an effect on expectation but increase variance.
 
\begin{itemize}
\item Episode 1: $\{(s,a),(s',a',R_1),(s,a,R_2),(s_t,R_3\}$ 
\begin{itemize}
\item First Visit
\bea
q(s,a) &=& R_1 + \gamma R_2 \frac{\pi(a'|s')}{b(a'|s')} +\gamma^2 R_3 \frac{\pi(a'|s') \pi(a|s)}{b(a'|s') b(a|s)} \\
q(s',a') &=& R_2 + \gamma R_3 \frac{\pi(a|s)}{b(a|s)}
\eea
\item Every Visit
\bea
q(s,a) &=& \frac{R_1 + \gamma R_2 \frac{\pi(a'|s')}{b(a'|s')} +\gamma^2 R_3 \frac{\pi(a'|s') \pi(a|s)}{b(a'|s') b(a|s)} +R_3}{2} \\
q(s',a') &=& R_2 + \gamma R_3 \frac{\pi(a|s)}{b(a|s)}
\eea
\end{itemize}
\item  Episode 2: $\{(s',a'),(s',a',R_4),(s,\bar a,R_5),(s_t,R_6)\}$
\begin{itemize}
\item First Visit
\bea
q(s,a) &=&R_1 + \gamma R_2 \frac{\pi(a'|s')}{b(a'|s')} +\gamma^2 R_3 \frac{\pi(a'|s') \pi(a|s)}{b(a'|s') b(a|s)} \\
q(s',a') &=& \frac{R_2 + \gamma R_3 \frac{\pi(a|s)}{b(a|s)} + R_4 + \gamma R_5 \frac{\pi(a'|s') }{b(a'|s') } + \gamma^2 R_6 \frac{\pi(a'|s') \pi(\bar a|s)}{b(a'|s') b(\bar a|s)}}{2} \\
q(s,\bar a) &=& R_6
\eea
\item Every Visit
\bea
q(s,a) &=&  \frac{R_1 + \gamma R_2 \frac{\pi(a'|s')}{b(a'|s')} +\gamma^2 R_3 \frac{\pi(a'|s') \pi(a|s)}{b(a'|s') b(a|s)} +R_3}{2} \\
q(s',a') &=& \frac{R_2 + \gamma R_3 \frac{\pi(a|s)}{b(a|s)} + R_4 + \gamma R_5 \frac{\pi(a'|s') }{b(a'|s') } + \gamma^2 R_6 \frac{\pi(a'|s') \pi(\bar a|s)}{b(a'|s') b(\bar a|s)} + R_5+ \gamma R_6 \frac{ \pi(\bar a|s)}{ b(\bar a|s)}}{3} \\
q(s,\bar a) &=& R_6
\eea
\end{itemize}
\end{itemize}


\subsection{Weighted Importance Sampling}

\bea
V(s) = \frac{\sum_{t \in \mathcal T(s)} \rho_{t:T(t)-1} G_t}{\sum_{t \in \mathcal T(s)} \rho_{t:T(t)-1} }
\eea
The way to interpret this is that $\mathcal T(s)$ is the union of all time steps at which the state was $s$.  If we are discussing first visit MC then this would be restricted to first visit. T(t) is the first time of termination following $t$. $G_t$ are returns pertaining to time step $t$ and $ \rho_{t:T(t)-1}$ are the importance sampling ratios.

To show how this works we consider an example. We show how the values are updated for first visit and every visit MC for the first two episodes. 
\begin{itemize}
\item Episode 1: $\{(s,a),(s',a',R_1),(s,a,R_2),(s_t,R_3\}$ 
\begin{itemize}
\item First Visit
\bea
q(s,a) &=& (R_1 + \gamma R_2+\gamma^2 R_3) \\
q(s',a') &=& (R_2 + \gamma R_3)
\eea
\item Every Visit
\bea
q(s,a) &=& \frac{(R_1 + \gamma R_2+\gamma^2 R_3) \frac{\pi(a'|s') \pi(a|s)}{b(a'|s') b(a|s)}+R_3}{\frac{\pi(a'|s') \pi(a|s)}{b(a'|s') b(a|s)}+1} \\
q(s',a') &=& (R_2 + \gamma R_3) 
\eea
\end{itemize}
\item  Episode 2: $\{(s',a'),(s',a',R_4),(s,\bar a,R_5),(s_t,R_6)\}$
\begin{itemize}
\item First Visit
\bea
q(s,a) &=& (R_1 + \gamma R_2+\gamma^2 R_3)  \\
q(s',a') &=& \frac{(R_2 + \gamma R_3) \frac{\pi(a|s)}{b(a|s)} + (R_4 + \gamma R_5 + \gamma^2 R_6) \frac{\pi(a'|s') \pi(\bar a|s)}{b(a'|s') b(\bar a|s)}}{ \frac{\pi(a|s)}{b(a|s)} + \frac{\pi(a'|s') \pi(\bar a|s)}{b(a'|s') b(\bar a|s)}} \\
q(s,\bar a) &=& R_6
\eea
\item Every Visit
\bea
q(s,a) &=& \frac{(R_1 + \gamma R_2+\gamma^2 R_3) \frac{\pi(a'|s') \pi(a|s)}{b(a'|s') b(a|s)}+R_3}{\frac{\pi(a'|s') \pi(a|s)}{b(a'|s') b(a|s)}+1} \\
q(s',a') &=& \frac{(R_2 + \gamma R_3) \frac{\pi(a|s)}{b(a|s)} + (R_4 + \gamma R_5 + \gamma^2 R_6) \frac{\pi(a'|s') \pi(\bar a|s)}{b(a'|s') b(\bar a|s)} + (R_5+ \gamma R_6)  \frac{ \pi(\bar a|s)}{ b(\bar a|s)}}{\frac{\pi(a|s)}{b(a|s)} + \frac{\pi(a'|s') \pi(\bar a|s)}{b(a'|s') b(\bar a|s)} +\frac{ \pi(\bar a|s)}{ b(\bar a|s)}} \\
q(s,\bar a) &=& R_6
\eea
\end{itemize}
\end{itemize}

\subsection{Algorithms}

\subsubsection{Every Visit Algorithms}
Here is the algorithm for on-policy MC prediction.

\begin{tcolorbox}[enhanced,title=On-policy prediction]
Initialize $Q(s,a)=$ random and $C(s,a)=0$ $\forall~s,a$ \\
Loop over episodes: \\
 \phantom{abcd} Using policy $\pi$ generate the sequence $S_0,A_0, R_1,S_1,A_1 \dots S_{T-1}, A_{T-1},R_T,S_T$ \\
 \phantom{abcd} $G=0$ \\
  \phantom{abcd} For $t$ in $(T-1,T-2,\dots 0)$: \\
  \phantom{abcd}   \phantom{abcd}  $G = \gamma G + R_{t+1}$ \\
  \phantom{abcd}   \phantom{abcd}  $Q(S_t,A_t) = \frac{C(S_t,A_t) Q(S_t,A_t) +G}{C(S_t,A_t)+1}$ \\
    \phantom{abcd}   \phantom{abcd}  $C(S_t,A_t) += 1$
\end{tcolorbox}

The algorithm for off-policy ordinary importance sampling MC prediction is

\begin{tcolorbox}[enhanced,title=Off-policy ordinary importance sampling prediction]
Initialize $Q(s,a)=$ random and $C(s,a)=0$ $\forall~s,a$ \\
Loop over episodes: \\
 \phantom{abcd} Using policy $b$ generate the sequence  $S_0,A_0, R_1,S_1,A_1 \dots S_{T-1}, A_{T-1},R_T,S_T$ \\
 \phantom{abcd} $G=0$ \\
  \phantom{abcd} $W=1$ \\
  \phantom{abcd} For $t$ in $(T-1,T-2,\dots 0)$ while $W \ne 0$: \\
  \phantom{abcd}   \phantom{abcd}  $G = \gamma G + R_{t+1}$ \\
  \phantom{abcd}   \phantom{abcd}  $Q(S_t,A_t) = \frac{C(S_t,A_t) Q(S_t,A_t) +W G}{C(S_t,A_t)+1}$ \\
    \phantom{abcd}   \phantom{abcd}  $C(S_t,A_t) += 1$ \\
    \phantom{abcd}   \phantom{abcd}  $W*= \frac{\pi(A_t|S_t)}{b(A_t|S_t)}$   
\end{tcolorbox}

The algorithm for off-policy weighted importance sampling MC prediction is

\begin{tcolorbox}[enhanced,title=Off-policy weighted importance sampling prediction]
Initialize $Q(s,a)=$ random and $C(s,a)=0$ $\forall~s,a$ \\
Loop over episodes: \\
 \phantom{abcd} Using policy $b$ generate the sequence  $S_0,A_0, R_1,S_1,A_1 \dots S_{T-1}, A_{T-1},R_T,S_T$ \\
 \phantom{abcd} $G=0$ \\
  \phantom{abcd} $W=1$ \\
  \phantom{abcd} For $t$ in $(T-1,T-2,\dots 0)$ while $W \ne 0$: \\
  \phantom{abcd}   \phantom{abcd}  $G = \gamma G + R_{t+1}$ \\
  \phantom{abcd}   \phantom{abcd}  $Q(S_t,A_t) = \frac{C(S_t,A_t) Q(S_t,A_t) +W G}{C(S_t,A_t)+W}$ \\
    \phantom{abcd}   \phantom{abcd}  $C(S_t,A_t) += W$ \\
    \phantom{abcd}   \phantom{abcd}  $W*= \frac{\pi(A_t|S_t)}{b(A_t|S_t)}$   
\end{tcolorbox}
The algorithm for off-policy per decision ordinary importance sampling MC prediction is

\begin{tcolorbox}[enhanced,title=Off-policy per decision ordinary importance sampling prediction]
Initialize $Q(s,a)=$ random and $C(s,a)=0$ $\forall~s,a$ \\
Loop over episodes: \\
 \phantom{abcd} Using policy $b$ generate the sequence  $S_0,A_0, R_1,S_1,A_1 \dots S_{T-1}, A_{T-1},R_T,S_T$ \\
 \phantom{abcd} $G=0$ \\
  \phantom{abcd} For $t$ in $(T-1,T-2,\dots 0)$: \\
  \phantom{abcd}   \phantom{abcd}  $G = \gamma G + R_{t+1}$ \\
  \phantom{abcd}   \phantom{abcd}  $Q(S_t,A_t) = \frac{C(S_t,A_t) Q(S_t,A_t) + G}{C(S_t,A_t)+1}$ \\
    \phantom{abcd}   \phantom{abcd}  $C(S_t,A_t) += 1$ \\
    \phantom{abcd}   \phantom{abcd}  $G*= \frac{\pi(A_t|S_t)}{b(A_t|S_t)}$   
\end{tcolorbox}

\subsection{Example}


Suppose we have a state $s$ and one terminal state $t$. There are only two actions left $L$ and right $R$. The left takes the agent to the terminal state with probability $q$ and return $1$. It takes the agent to the state $s$ with probability $1-q$. Suppose further the target policy has $\pi(L|s)=p$ and the exploration policy has $b(L|s)=\tilde p$. We evaluate the value and variance of the state under the first visit method. Then we have

\bea
\mathbb E_\pi[G_0] &=& pq \sum_{k=0}^\infty (p (1-q))^k \nn
&=& \frac{pq}{1-p(1-q)}
\eea
We also have
\bea
\mathbb E_\pi[G_0^2] &=& pq \sum_{k=0}^\infty (p (1-q))^k \nn
&=& \frac{pq}{1-p(1-q)}
\eea

If we compute this with importance sampling we get
\bea
\mathbb E_b[G_0 \prod_{t=0}^\infty \frac{\pi(A_t|S_t)}{b(A_t|S_t)}] &=& \tilde pq \sum_{k=0}^\infty (\tilde p (1-q))^k \left( \frac{p}{\tilde p} \right)^{k+1} \nn
&=& \frac{pq}{1-p(1-q)}
\eea
which is the same as under the target policy.

We also get
\bea
\mathbb E_b[\left(G_0 \prod_{t=0}^\infty \frac{\pi(A_t|S_t)}{b(A_t|S_t)}\right)^2] &=& \tilde pq \sum_{k=0}^\infty (\tilde p (1-q))^k \left( \frac{p}{\tilde p} \right)^{2(k+1)} \nn
&=& \frac{qp^2}{\tilde p} \sum_{k=0}^\infty \left( \frac{(1-q) p^2}{\tilde p}\right)^k
\eea
This variance is divergent if $(1-q)p^2 \ge \tilde p$. In Barto and Sutton, they take $q=0.1, p=1$ and $\tilde p=0.5$.


\section{Policy Gradient}


Below we understand that $\pi(a|s)$ depends on $\theta$ parameterically.

We then have
\bea
\partial_\theta v_\pi(s) &=& \partial_\theta \sum_a q_\pi(s,a) \pi(a|s) \nn
&=& \sum_a \partial_\theta \pi(a|s) q_\pi(s,a) + \sum_a \pi(a|s) \partial_\theta q_\pi(s,a) \nn
&=& \sum_a \partial_\theta \pi(a|s) q_\pi(s,a)  + \gamma \sum_{a,s'} \pi(a|s) p(s'|s,a) \partial_\theta v_\pi(s') \nn
&=& \sum_a \partial_\theta \pi(a|s) q_\pi(s,a) + \gamma \sum_{s'} \partial_\theta v_\pi(s') p(s'|s)
\eea
Thus, iterating the recursion and noting that a state $x$ reached from $s$ in $k$ steps accumulates a factor $\gamma^k$, we define the (unnormalised) discounted visitation $p(x|s) \coloneqq \sum_{k=0}^\infty \gamma^k p_k(x|s)$, where $p_k(x|s)$ is the $k$-step transition probability under $\pi$. This gives
\bea
\partial_\theta v_\pi(s) &=& \sum_{x \in \mathcal S} p(x|s) \sum_a \partial_\theta \pi(a| x) q_\pi(x,a) \nn
&=& \mathbb E_\pi \left[  \sum_a \partial_\theta \pi(a| x) q_\pi(x,a) \right] \nn
&=& \mathbb E_\pi \left[ \sum_a \pi(a| x) \partial_\theta \left( \log \pi(a|x) \right)q_\pi(x,a) \right] \nn
&=& \mathbb E_\pi \left[ \mathbb E_a [ \partial_\theta \left( \log \pi(a|x) \right)q_\pi(x,a)]  \right] \nn
\eea





\section{Importance Sampling Example}
Note that
\bea
v_\pi(s) = \mathbb E_b [ \rho_{t:T-1} G_t | S_t=s]
\eea
where $\rho_{t:T-1} = \prod_{k=t}^{T-1} \frac{\pi(A_k|S_k)}{b(A_k|S_k)}$ and
\bea
q_\pi(s,a) = \mathbb E_b [ \rho_{t+1:T-1} G_t | S_t=s, A_t=a]
\eea


Suppose we have a sequence $S_0,A_0,R_1,S_1,A_1,R_2,S_2,A_2,R_3,S_3$ where $S_0=S_2=s$ and $A_0=A_2=a$ and $S_1=s',A_1=a'$ then for ordinary importance sampling starting with $Q(S,A)=0$ we get
\bea
q(s,a) &=& \frac{G_2 + G_0 \rho_{1:2}}{2} \nn
&=&  \frac{ R_3 + (R_1 + \gamma R_2 + \gamma^2 R_3) \frac{\pi(A_1|S_1)\pi(A_2|S_2)}{b(A_1|S_1)b(A_2|S_2)}}{2} \\
q(s',a') &=& G_1 \rho_{2:2} \nn
&=& (R_2 + \gamma R_3) \frac{\pi(A_2|S_2)}{b(A_2|S_2)}
\eea
whereas for weighted importance sampling we have

\bea
q(s,a) &=&  \frac{G_2 + G_0 \rho_{1:2}}{1+\rho_{1:2}} \nn
&=&\frac{ R_3 + (R_1 + \gamma R_2 + \gamma^2 R_3) \frac{\pi(A_1|S_1)\pi(A_2|S_2)}{b(A_1|S_1)b(A_2|S_2)}}{1+\frac{\pi(A_1|S_1)\pi(A_2|S_2)}{b(A_1|S_1)b(A_2|S_2)}} \\
q(s',a') &=& G_1 \nn
&=& (R_2 + \gamma R_3)
\eea

Now note that because of causality we have
\bea
\mathbb E_b [ \rho_{t:T-1}R_{t+k} ] = \mathbb E_b [ \rho_{t:t+k-1}R_{t+k} ]
\eea
so we could also simply take

\bea
q(s,a) &=&  \frac{ R_3 + (R_1 + \gamma R_2 \frac{\pi(A_1|S_1)}{b(A_1|S_1)} + \gamma^2 R_3 \frac{\pi(A_1|S_1)\pi(A_2|S_2)}{b(A_1|S_1)b(A_2|S_2)} ) }{2} \\
q(s',a') &=& (R_2 + \gamma R_3 \frac{\pi(A_2|S_2)}{b(A_2|S_2)})
\eea



\end{document}